% s-374-p374-m001; original PDF 374
\PoemStart{s-374-p374-m001}{Alpenleben.}
\begin{BookVerse}
\VerseLine{0}{Der Mittag hat der Hölle Gluth entriegelt,}
\VerseLine{1}{Die Wiese ist ein fieberheisser Pfühl;}
\VerseLine{1}{Der Himmel trieft vor Bläue, lodert schwül,}
\VerseLine{0}{Vom Rund der Sonne feuerfest versiegelt.}
\Stanza
\VerseLine{0}{Doch war sie doppelt wild und aufgewiegelt,}
\VerseLine{0}{Wird Abends doppelt still die Welt und kühl;}
\VerseLine{1}{Die Unruh’ glättet sich, und das Gefühl}
\VerseLine{0}{Beschaut im See des Friedens sich gespiegelt.}
\Stanza
\VerseLine{0}{Thalab geschlendert kommen die Gedanken,}
\VerseLine{1}{Gleich Kühen satt und ohne Leidenschaften.}
\VerseLine{0}{Mit Lust seh’ ich die vollen Euter schwanken.}
\Stanza
\VerseLine{0}{Was kümmert’s mich, in welch besonnten Räumen}
\VerseLine{1}{Sie Mittags sich ihr Weideglück errafften?}
\VerseLine{0}{Mir soll die kühle Milch im Kruge schäumen!}
\end{BookVerse}
\PoemEnd
